\section{Data and Measurement}\label{sec:data}

\paragraph{CPS monthly couples.} The CPS basic monthly files for 1994--2023 identify both spouses of a married couple
and follow them for four consecutive months, and again for four months a year later. I link married couples with both
spouses present and aged 20--60 from month to month within the rotation and across the twelve-month break, which gives
6.7 million couple-months of 1.3 million couples. A spouse loses a job in month $t$ if he or she was employed in $t-1$
and is unemployed in $t$ after a layoff, a job loss or the end of a temporary job; quits and exits from the labor force
are not counted as job losses and are reported separately. The monthly design measures timing to the month and allows
state $\times$ month and state $\times$ month $\times$ industry fixed effects. Its limits are the four-month window, the
absence of an employer identifier and proxy reporting: one household member usually answers for both spouses, so
reporting errors could be correlated, and part of the job losses are temporary layoffs. Section~\ref{sec:jointloss}
shows that neither drives the result.\footnote{Linear probability models; standard errors clustered by household (clustering by state
gives the same significance). Variables, link rates and the months in which household identifiers break are documented
in \texttt{IGE/output/cps\_joint\_job\_loss.md}. Code: \texttt{IGE/30.CPSJointJobLoss.py}.}

\paragraph{CPS Displaced Worker Supplement.} The biennial supplement (1984--2024) records, for each adult, whether a job
was lost in the last three years because the plant closed or moved, there was insufficient work or the position was
abolished, the year of the loss and the industry of the lost job. Both spouses answer it, so it measures joint
displacements with each spouse's own reason in a sample of 417,269 couples.

\paragraph{SIPP.} The 2004, 2008 and 2014 panels of the Survey of Income and Program Participation follow both
spouses of 50,692 married couples month by month for up to five years, with the reason each job ended. They measure the
joint separation to the month over a longer window than the CPS and trace the wife's response for three years.

\paragraph{PSID displacements.} The PSID asks household heads, and from 1984 their wives, why their last job ended. I
code a displacement when the job ended because the company closed or moved, or because the person was laid off or
fired, and date it to the calendar year in which the job ended. Closures alone serve as the strictest definition and
quits as a placebo. Among working male heads aged 20--60, the annual displacement rate is 8.0\% for Black and 4.8\%
for White men.\footnote{The questions differ across waves: before 1984 they refer to the previous job of heads who
started their current job in the last year and to the last job of unemployed heads; from 1984 to 2001 they cover
previous employers and positions with end years; from 2003 they cover each employer in the employment history. Three
quarters of the reports carry the end year; the rest are dated to the year before the interview. Code:
\texttt{IGE/19.FatherDisplacement.py}.} The PSID follows couples for decades, so it measures the household's response
over years and the children's adult outcomes, with the SRC and SEO samples providing Black and White families.\footnote{The
National Longitudinal Surveys of Youth (NLSY79, NLSY97) are not used for the joint job loss. They record the
respondent's own jobs in great detail, including weekly employment, an employer roster and the reason each job ended,
but only a few annual items on the spouse (whether the spouse worked last week, weeks worked and earnings last year),
with no spouse employer roster and no reason or date for the spouse's job losses. They therefore cannot measure whether
both spouses lost their jobs, when and for which reason, which the PSID, the CPS, the Displaced Worker Supplement and
the SIPP measure for both spouses.}

\paragraph{Event studies.} For the household I estimate a stacked event study of the father's first displacement: for
each displaced father with a child aged 0--17 at home, the controls are fathers not displaced at any time up to six
years after the event year; event time runs from $-3$ to $+6$, relative to year $-2$. For the joint job loss the unit is
a father $\times$ event year with his wife employed at baseline, and the outcome is the wife's own displacement in the
three years around the event, with couple and calendar-year fixed effects in the event-study version and a placebo
window before the event.

\paragraph{Intergenerational persistence.} The children's adult outcomes are linked to their parents' income: in the PSID the IGE of sons' family income at 40 is 0.27 for Black
and 0.47 for White sons, and at the 25th percentile of parental income Black children reach a household income rank
12.9 percentiles below White children, against 12.6 in \citet{chetty2020race}.\footnote{Code:
\texttt{IGE/6.IGEMain.py} and \texttt{IGE/7.RankMobility.py}; details in Appendix~\ref{app:other}.}
